\begin{frame}{미로 탐색의 Mission}\textbf{Input:} grid, start $(r,c)$, exit.\quad \textbf{Output:} 경로 존재 여부.\mission{\texttt{path(r,c)}는 현재 cell에서 출구로 갈 수 있으면 true를 반환한다.}
\end{frame}
\begin{frame}{미로 경로의 재귀적 정의}경로가 존재하려면 현재 cell이 출구이거나, 방문하지 않은 통로 이웃 중 하나에서 출구까지 경로가 존재한다.\[P(r,c)=exit(r,c)\lor\bigvee_{(r',c')\in N(r,c)}P(r',c').\]
\end{frame}
\begin{frame}{왜 visited가 필요한가?}격자에는 cycle이 있다: A→B→A→B… .\begin{alertblock}{규칙}재귀 호출 \textbf{전에} 현재 cell을 visited로 표시한다.\end{alertblock}그러면 각 cell은 최대 한 번 확장된다.
\end{frame}
\begin{frame}[fragile]{경로 존재 여부: Pseudocode}\begin{lstlisting}
bool find_path(int r, int c) {
  if (outside(r,c) || wall(r,c) || visited(r,c)) return false;
  if (is_exit(r,c)) return true;
  mark_visited(r,c);
  for each (nr,nc) in four_neighbors(r,c)
    if (find_path(nr,nc)) return true;
  return false;
}
\end{lstlisting}
\vfill\muted{Boolean 존재 여부만 반환하므로 visited를 unmark할 필요가 없다. 경로 출력용 current path와는 역할이 다르다.}
\end{frame}
\begin{frame}{경로를 남기려면}\begin{enumerate}\item current를 \textbf{candidate path}로 표시\item 이웃을 재귀적으로 explore\item 모두 실패하면 \textbf{dead end}로 바꾸고 돌아감\end{enumerate}성공 상태와 방문했지만 실패한 상태를 구분한다.
\end{frame}
\begin{frame}{6×6 미로: 탐색과 되돌아가기}
\begin{columns}[c,onlytextwidth]
\column{.62\textwidth}\centering
\begin{tikzpicture}[scale=1.08,transform shape]
\foreach \r in {0,...,5}{\foreach \c in {0,...,5}{\gridcell{\r}{\c}{open}}}
\foreach \r/\c in {0/2,0/3,0/4,0/5,1/0,1/2,1/3,1/4,1/5,2/4,2/5,3/0,3/1,3/4,3/5,4/0,4/1,4/5,5/0,5/1,5/2,5/3}{\gridcell{\r}{\c}{wall}}
\only<1|handout:0>{\foreach \r/\c in {0/0,0/1,1/1}{\gridcell{\r}{\c}{path}}}
\only<2|handout:0>{\foreach \r/\c in {0/0,0/1,1/1,2/1}{\gridcell{\r}{\c}{path}}\gridcell{2}{0}{dead}}
\only<3|handout:0>{\foreach \r/\c in {0/0,0/1,1/1,2/1,2/2,3/2,4/2,4/3}{\gridcell{\r}{\c}{path}}\gridcell{2}{0}{dead}}
\only<4|handout:0>{\foreach \r/\c in {0/0,0/1,1/1,2/1,2/2,3/2,4/2,4/3,3/3,2/3}{\gridcell{\r}{\c}{path}}\gridcell{2}{0}{dead}}
\only<5|handout:0>{\foreach \r/\c in {0/0,0/1,1/1,2/1,2/2,3/2,4/2,4/3}{\gridcell{\r}{\c}{path}}\foreach \r/\c in {2/0,2/3,3/3}{\gridcell{\r}{\c}{dead}}}
\only<6|handout:1>{\foreach \r/\c in {0/0,0/1,1/1,2/1,2/2,3/2,4/2,4/3,4/4,5/4,5/5}{\gridcell{\r}{\c}{path}}\foreach \r/\c in {2/0,2/3,3/3}{\gridcell{\r}{\c}{dead}}}
\node[font=\bfseries] at (0,0){S};\node[font=\bfseries] at (5,-5){E};
\end{tikzpicture}
\column{.34\textwidth}
\only<1|handout:0>{\Large 후보 경로를 따라 전진}
\only<2|handout:0>{\Large 막힌 가지를 \textbf{dead end}로 표시}
\only<3|handout:0>{\Large 아래로 내려가 오른쪽 통로로 꺾어 전진}
\only<4|handout:0>{\Large 위쪽 통로로 전진을 시도하는 중}
\only<5|handout:0>{\Large 더 나아갈 수 없어 되돌아가며 지나온 칸을 dead로 표시}
\only<6|handout:1>{\Large 출구까지 한 경로를 발견}
\end{columns}
\end{frame}
\begin{frame}{Backtracking 패턴}\[\boxed{\text{choose}\;\to\;\text{explore}\;\to\;\text{unchoose/mark dead}}\]선택이 실패하면 호출 전 상태로 복구하거나, 다시 탐색할 필요가 없도록 blocked로 확정한다. 미로 탐색은 grid 위의 DFS다.
\end{frame}
\begin{frame}{미로 탐색의 복잡도}\begin{itemize}\item 상하좌우의 constant-degree adjacency를 사용한다.\item visited를 사용하면 각 열린 cell을 최대 한 번 확장: $O(RC)$\item 저장 grid 외 recursion stack: 최악 $O(RC)$\item visited가 없으면 종료하지 않거나 같은 경로를 반복\end{itemize}\sectiontakeaway{정확한 상태 표현이 종료성과 효율성을 동시에 만든다.}
\vfill\muted{경로 존재 여부를 찾는 DFS를, 연결된 모든 cell을 세는 DFS로 확장한다.}
\end{frame}
